\chapter{\texttt{core.perturbative} --- perturbative approximation}
\label{ch:core-perturbative}

\begin{notebox}[title={This Chapter Describes an Approximation}]
Everything in this chapter is a \textbf{first-order perturbative
approximation} to the exact position-dependent Schr\"odinger equation
(\cref{ch:core-common}), offered as a cheaper alternative to the exact
numerical method of \cref{ch:core-numerical} \emph{within a restricted
validity range} -- it is not a second, independently correct propagation
method.

\textbf{The approximation.} Within one segment, the Hamiltonian is split
into a constant average and a position-dependent residual,
$H(x)=H_{\rm avg}+\delta H(x)$, where $H_{\rm avg}$ uses the
segment-\emph{average} density and $\delta H(x)$ captures the fluctuation
around it. The evolutor is expanded to \textbf{first order} in $\delta H$:
$S\approx U_0+U_1$, with $U_0=\exp(-iH_{\rm avg}L)$ exact for the constant
average (\cref{sec:perturbative-evolutor}), and $U_1$ the first Dyson-series
correction from $\delta H$.

\textbf{Range and conditions of validity.} The truncation is controlled by
the phase the \emph{residual} potential accumulates over the segment,
$|\delta V(x)|\cdot L$: the approximation is trustworthy while this stays
small compared to 1 throughout the segment (equivalently, while the
neglected second- and higher-order Dyson terms stay small relative to
$U_1$). In practice this means: (a) the profile model
(\cref{sec:perturbative-models}) must fit the true density well enough,
and/or segments must be short enough, that the density does not deviate
strongly from its local average within any one segment; (b) accuracy is
expected to degrade for segments spanning a large density contrast (e.g.\ a
sharp discontinuity treated as a single segment). \textbf{No automatic
error estimate or order check is built into this code path} -- unlike
\cref{ch:core-numerical}'s explicit accuracy note, there is no signal here
that a given (segment, profile-model) choice has left the valid regime; it
must be established externally, typically by cross-validating against the
exact numerical evolutor of \cref{ch:core-numerical} for representative
parameters before trusting a perturbative configuration in production.
\end{notebox}

This chapter covers \texttt{spectral.py} (the closed-form spectral
decomposition every perturbative segment evolutor is built from),
\texttt{evolutor.py} (the zeroth- and first-order evolution operators
themselves), and the family of interchangeable density-profile models that
supply the analytic inputs the first-order correction needs.

\section{\texttt{spectral.py} --- closed-form spectral decomposition}

\modulehead{core/perturbative/spectral.py}{Trace/traceless splitting,
characteristic-polynomial invariants, and closed-form spectral projectors
for a $3\times3$ or $4\times4$ Hermitian Hamiltonian.}

\begin{theorybox}
\subsection*{Trace and traceless parts}
Any Hermitian $H$ splits uniquely into a multiple of the identity and a
traceless remainder,
\begin{equation}
  \keyresult{T = H - \frac{\Tr(H)}{N}\,I_N}, \qquad N=H.\text{shape}[-1]\in\{3,4\}.
\end{equation}
Since $I_N$ commutes with everything, $H$ and $T$ share the same
eigenvectors, and $H$'s eigenvalues are simply $T$'s shifted uniformly by
$\Tr(H)/N$. Working with $T$ instead of $H$ removes the sub-leading term of
the characteristic polynomial (no $\lambda^{N-1}$ term, since $\Tr(T)=0$),
which is what makes the closed-form projector formulas below tractable.

\subsection*{Characteristic-polynomial invariants}
For $N=3$, $T$'s characteristic polynomial is $\lambda^3+c_1\lambda+c_0=0$
with
\begin{equation}
  c_1 = -\tfrac12\Tr(T^2), \qquad c_0 = -\tfrac13\Tr(T^3).
\end{equation}
The $N=4$ projector formula additionally needs
$e_3=\tfrac13\Tr(T^3)=-c_0$ (kept as a separate function in the code to
avoid a sign-confusion bug, since both invariants are used together for
$N=4$). The eigenvalues $\lambda_a$ themselves have no closed-form solver
in this module -- they are obtained numerically via
\pyfunc{torch.linalg.eigvalsh}; only the \emph{projectors} below are
closed-form, given the eigenvalues.

\subsection*{Closed-form spectral projectors}
The spectral theorem writes $T=\sum_a\lambda_a M_a$ with
$M_a=\prod_{b\neq a}\frac{T-\lambda_b I}{\lambda_a-\lambda_b}$ (the
Lagrange/Sylvester interpolation projector). Using Cayley--Hamilton and
$\Tr(T)=0$ to eliminate the "other" eigenvalues $\lambda_{b\neq a}$ from
this product, it reduces to an explicit polynomial in $T$ (and $T^2$, plus
$T^3$ for $N=4$) with coefficients depending only on $\lambda_a$ and the
invariants above:
For $N=3$:
\begin{equation}
  \keyresult{M_a = \frac{(\lambda_a^2+c_1)\,I + \lambda_a T + T^2}{3\lambda_a^2+c_1}}.
\end{equation}
For $N=4$:
\begin{equation}
  \keyresult{M_a = \frac{T^3+\lambda_a T^2+(c_1+\lambda_a^2)T+(\lambda_a^3+c_1\lambda_a-e_3)\,I}{4\lambda_a^3+2c_1\lambda_a-e_3}}.
\end{equation}
Both denominators equal $p'(\lambda_a)$, the derivative of the
characteristic polynomial at that root: they vanish exactly when
$\lambda_a$ is a repeated eigenvalue. Both formulas were verified
symbolically and numerically (against direct Lagrange interpolation and
against \pyfunc{torch.linalg.eigh} eigenvector outer products) to match to
floating-point precision.
\end{theorybox}

\begin{notebox}
\textbf{Degeneracy fallback.} Near a (near-)degenerate spectrum the
closed-form projector formulas lose precision well before their
denominator literally reaches zero, since $p'(\lambda_a)\to0$
continuously. \pyfunc{\_spectral\_degeneracy\_mask} flags batch entries
whose minimum pairwise eigenvalue gap is small \emph{relative} to the
spectrum scale (or whose scale itself is numerically zero, the $T\approx0$
case, where the relative test would be $0/0$); flagged entries have their
projectors replaced wholesale by $M_a=v_av_a^\dagger$ from a direct
\pyfunc{torch.linalg.eigh(T)} call -- exact and stable by construction,
independent of the closed-form formula's conditioning, computed only when
at least one batch entry actually needs it. This module runs entirely
under \code{@torch.no\_grad()}, so the well-known gradient instability of
degenerate eigenvectors is not a concern here (the perturbative path is
used for forward propagation, not for differentiating through the
spectrum).
\end{notebox}

\begin{implbox}
\begin{itemize}
  \item \pyfunc{hamiltonian\_traceless(H)}: returns $(T,\Tr H)$.
  \item \pyfunc{hamiltonian\_traceless\_c0/c1/e3(T)}: the invariants above.
  \item \pyfunc{hamiltonian\_traceless\_eigenvalues(T)}: $\lambda_a$ via
    \pyfunc{torch.linalg.eigvalsh}.
  \item \pyfunc{hamiltonian\_spectral\_projectors\_traceless(T, ...)}:
    builds $M_a$ (closed form, with the \pyfunc{eigh} fallback above);
    returns $(M,\lambda,c_1)$.
  \item \pyfunc{hamiltonian\_spectral\_data(H)}: the one-stop entry point
    used by \texttt{evolutor.py} -- computes $T$, $\Tr H$, $\lambda$, $c_1$,
    and $M$ in one call, reusing intermediate products ($T^2$, $T^3$)
    across every invariant and the projector formula to avoid redundant
    matrix multiplications.
  \item \pyfunc{spectral\_projector\_residuals(M, T, lam)}: diagnostic
    residual norms (completeness $\sum_aM_a=I$, idempotency
    $M_a^2=M_a$, orthogonality $M_aM_b=0$, unit trace, and
    reconstruction $\sum_a\lambda_aM_a=T$) -- used both by the test suite
    and by the degeneracy instrumentation itself.
\end{itemize}
\end{implbox}

\section{\texttt{evolutor.py} --- perturbative segment evolution}
\label{sec:perturbative-evolutor}

\modulehead{core/perturbative/evolutor.py}{\pyfunc{evolutor\_zero\_order}
(exact, for the segment-average Hamiltonian) and
\pyfunc{evolutor\_first\_order} (the first Dyson-series correction from the
density residual), combined into the segment evolutor below.}

\begin{theorybox}
\subsection*{Zeroth order: exact evolution under the average Hamiltonian}
Given the spectral decomposition $H_{\rm avg}=T+\Tr(H_{\rm avg})/N\cdot I$
from the previous section, the spectral theorem gives $\exp(-iH_{\rm avg}L)$
\emph{exactly} (not perturbatively) as
\begin{equation}
  \keyresult{U_0 = \sum_a \exp\!\Bigl[-i\Bigl(\lambda_a+\tfrac{\Tr H_{\rm avg}}{N}\Bigr)L\Bigr]\,M_a},
\end{equation}
since conjugating by the identity-shift only shifts every eigenvalue by the
same constant $\Tr(H_{\rm avg})/N$, leaving the projectors $M_a$
untouched. $U_0$ is the entire answer if the segment's density were exactly
constant at its average value; every approximation in this chapter lives
in the correction below.

\subsection*{First order: the Dyson-series correction from the residual}
Writing $H(x)=H_{\rm avg}+\delta H(x)$ with $\delta H(x)=\delta V(x)\,P$ (a
position-independent \emph{direction} matrix $P$ scaled by the
position-dependent residual potential $\delta V(x)=V(x)-V_{\rm avg}$; see
below for $P$), first-order time-dependent perturbation theory in the
interaction picture defined by $U_0$ gives
\begin{equation}
  U_1 = -i\int_{x_1}^{x_2} U_0(x_2,x')\,\delta H(x')\,U_0(x',x_1)\,\mathrm{d}x'.
\end{equation}
Expanding $U_0$ in the spectral basis, $U_0(x_2,x')=\sum_a
e^{-i\lambda_a(x_2-x')}M_a$, and inserting $\delta H(x')=\delta V(x')P$
turns the operator integral into a sum over eigenvalue \emph{pairs} of a
scalar oscillatory integral of the residual density weighted by the matrix
$M_aPM_b$:
\begin{equation}
  \keyresult{U_1 = -i\sum_{a,b} I_{ab}\;M_a\,P\,M_b},
  \qquad
  I_{ab} = \int_{x_1}^{x_2}\delta V(x')\,e^{-i\lambda_a(x_2-x')-i\lambda_b(x'-x_1)}\,\mathrm{d}x'.
\end{equation}
$I_{ab}$ (a ``spectral oscillatory integral'' of the residual profile) is
computed analytically by the active profile model
(\cref{sec:perturbative-models}, \pyfunc{residual\_integral}); it is the
only place the specific density parametrisation enters the perturbative
machinery -- \texttt{evolutor.py} itself never sees a profile's functional
form, only $I_{ab}(\lambda_a,\lambda_b)$.

\textbf{The direction matrix $P$.} For the plain Standard Model or 3+1
sterile case (no NSI), only the electron-flavour entry of $H_{\rm mat}$
varies with density, so $P=\diag(1,0,\dots,0)$ and the sandwich $M_aPM_b$
reduces to the rank-1 outer product of $M_a$'s first column with $M_b$'s
first row -- the closed-form shortcut used when the caller passes no
explicit $P$. When NSI is active, the coupling $\epsilon$ is
position-independent (only the electron density magnitude, not the NSI
coupling itself, varies along the segment), so the same first-order
structure applies with
\begin{equation}
  P = O^\dagger\bigl(\diag(1,0,\dots,0)+\epsilon\bigr)O
\end{equation}
in the reduced basis (\cref{ch:core-bsm}) -- fixed once per segment (built
at $V{=}1$) and reused for both $H_{\rm avg}$'s matter term and this
first-order sandwich, since it no longer reduces to the SM rank-1
shortcut once $\epsilon$ breaks the electron-only structure.

\textbf{The neutral-current direction $P_{NC}$.} When the optional 3+1
sterile neutral-current term is active (\cref{ch:core-common}), it
contributes a second, wholly independent first-order term, since the
electron- and neutron-density residuals fluctuate independently along the
segment:
\begin{equation}
  \keyresult{U_1 = -i\sum_{a,b} I_{ab}^{(e)}\;M_a\,P_{CC}\,M_b
             \;-i\sum_{a,b} I_{ab}^{(n)}\;M_a\,P_{NC}\,M_b},
\end{equation}
where $I_{ab}^{(n)}$ is the same spectral oscillatory integral as
$I_{ab}^{(e)}$ but of the \emph{neutron}-density residual
($\pyfunc{residual\_integral\_neutron}$, below), and
$P_{NC}=O^\dagger\diag(0,0,0,-1)\,O$ is the (always rotated) direction
matrix derived in \cref{ch:core-common}. Unlike $P_{CC}$, $P_{NC}$ never
reduces to a rank-1 shortcut regardless of whether NSI is active, since it
never commutes with $O$; the CC term above may still use its own rank-1
shortcut independently, since the two contributions are additive and
computed separately. Omitting the neutral-current term (the default)
recovers the CC(+NSI)-only expression exactly.

\subsection*{Combining orders, and the zero-length convention}
The segment evolutor is $U=U_0+U_1$
(\pyfunc{evolutor\_perturbative\_from\_H}), with two guards applied by the
code around this core formula: (1) if the profile model reports no
perturbation for a segment (a genuinely constant density, e.g.\
\pyfunc{has\_perturbation()} is \code{False}), $U_1$ is skipped entirely
--- both because it is identically zero and to avoid computing
$I_{ab}$ needlessly; (2) zero-length segments are forced to the identity
operator (\pyfunc{enforce\_identity\_for\_zero\_length}) regardless of the
formula above, since $L=0$ is a genuine edge case (e.g.\ a trajectory that
grazes but does not cross a layer) rather than a regime where the
perturbative expansion itself is being tested.
\end{theorybox}

\begin{notebox}
\textbf{Hermitian conjugate, again.} \pyfunc{evolutor\_perturbative\_segment}
builds $H_{\rm kin}$ via \pyfunc{hamiltonian\_kinetic\_reduced}
(\cref{ch:core-common}), which -- exactly as flagged there and in
\cref{ch:core-bsm} -- always uses $U_{\rm red}\diag(k)U_{\rm
red}^\dagger$, never the plain transpose. This matters even more visibly
here: for the 3+1 sterile case the code computes the segment's flavour-basis
trace directly as $\Tr(H_{\rm kin}+H_{\rm mat})$ rather than assuming
$\Tr(H_{\rm kin})=\sum_ik_i$ (only guaranteed for a \emph{real} $U_{\rm
red}$) or $\Tr(H_{\rm mat})=V$ (only true for the plain
$\diag(1,0,\dots,0)$ direction) -- both shortcuts used safely in the
pure-SM fast path, both invalid once the sterile CP phase $\delta_{14}$ or
a generic NSI $\epsilon$ is active.
\end{notebox}

\begin{implbox}
\begin{itemize}
  \item \pyfunc{evolutor\_zero\_order(H, L, ...)}: the exact $U_0$ formula
    above; can also return the spectral data for reuse by the first-order
    call, avoiding a second diagonalisation.
  \item \pyfunc{evolutor\_first\_order(M, lam, trace\_H, profile\_model, ..., P=None, P\_nc=None)}:
    the $U_1$ formula above, using the rank-1 shortcut for the CC term when
    \code{P} is \code{None}; when \code{P\_nc} is also supplied, adds the
    independent neutral-current sandwich term using
    \pyfunc{profile\_model.residual\_integral\_neutron} and
    \pyfunc{matter\_potential\_nc}.
  \item \pyfunc{evolutor\_perturbative\_from\_H(H, L, profile\_model, ..., P=None, P\_nc=None)}:
    combines both orders, with the no-perturbation and zero-length guards;
    the no-perturbation check also consults
    \pyfunc{profile\_model.has\_perturbation\_neutron()} when \code{P\_nc}
    is supplied.
  \item \pyfunc{evolutor\_perturbative\_segment(oscillation, E\_MeV,
    profile\_model, ..., include\_matter\_nc=False)}: the user-facing entry
    point -- builds the reduced Hamiltonian (SM, sterile, NSI, or any
    combination, mirroring \cref{ch:core-common,ch:core-bsm}) from the
    profile model's segment average, then evolves it perturbatively via the
    function above. The plain Standard Model case takes a dedicated fast
    path (bit-identical output, avoiding the general-$N$ machinery's
    overhead). \code{include\_matter\_nc=True} additionally requires
    \code{profile\_model.potential\_n} to be set (i.e.\ the profile was
    built with neutron-density coefficients, next section) and is silently
    ignored for a 3-flavour \code{oscillation.pmns}, mirroring
    \pyfunc{hamiltonian\_matter\_reduced}'s own convention
    (\cref{ch:core-common}).
\end{itemize}
\end{implbox}

\section{Comparing the perturbative density-profile models}
\label{sec:perturbative-models}

\modulehead{core/perturbative/models/*}{Three interchangeable profile
models, each supplying \code{average}, \code{length}, \code{zero\_mask},
\code{potential}, \pyfunc{residual\_integral}, and
\pyfunc{has\_perturbation} to the model-agnostic evolutor above.}

\begin{implbox}
None of the mathematics in \texttt{evolutor.py} or \texttt{spectral.py}
changes between models: they differ purely in \emph{how} the segment's
density profile $n_e(x)$ is parametrised, and hence in the closed-form
expression used for $I_{ab}$.

\begin{center}
\small
\begin{tabular}{@{}ll@{}}
\toprule
Model & Density parametrisation and usage \\
\midrule
\parbox[t]{3.2cm}{\code{even\_power}} &
  \parbox[t]{9.4cm}{$n_e(x)=a+bx^2+cx^4+\delta_1x^6+\dots$ (even powers only,
  arbitrary order) in the raw trajectory coordinate $x$. Appropriate
  whenever the segment is symmetric about the trajectory's point of
  closest approach to a spherically symmetric medium (e.g.\ an
  Earth-crossing chord through one shell), which forces $n_e(x)=n_e(-x)$
  and eliminates odd powers identically. The general-order model used for
  Earth propagation.} \\[30pt]
\parbox[t]{3.2cm}{\code{prem}} &
  \parbox[t]{9.4cm}{$n_e(x)=a+bx^2$: the 2-coefficient special case of
  \code{even\_power} truncated at the quadratic term, matching the PREM
  reference model's own native linear-in-$r^2$ parametrisation per shell
  \parencite{DziewonskiAnderson1981}. A lighter-weight specialisation, not
  an independent formulation -- it exists to avoid the general
  \code{even\_power} machinery's overhead when only two coefficients are
  ever needed.} \\[30pt]
\parbox[t]{3.2cm}{\code{atmosphere}} &
  \parbox[t]{9.4cm}{$n_e(q)=\sum_kc_kq^k$ (every power, even and odd) in a
  \emph{centred and rescaled} local coordinate $q=(x-\text{centre})/\text{half-width}$.
  Atmosphere segments are not symmetric about their own centre the way an
  Earth chord is, so odd terms are required; the local coordinate keeps the
  fitted coefficients well-conditioned regardless of the segment's absolute
  position or width.} \\[30pt]
\bottomrule
\end{tabular}
\end{center}

Every model computes its own \pyfunc{residual\_integral} in closed form for
ordinary (non-degenerate) eigenvalue gaps, switching to a Taylor expansion
in $\lambda_a-\lambda_b$ near degeneracy to avoid the numerical instability
of the closed form as the oscillation frequency approaches zero -- the same
degeneracy-awareness pattern as the spectral projectors themselves
(previous section), applied independently at the profile-model level.
\pyfunc{perturbative\_profile\_selection} (\code{core/perturbative/models/model\_selection.py})
maps a public model name (\code{"even\_power"}, \code{"prem"}, ...) to the
concrete layered-profile class, so medium code can request a model without
importing model-specific classes directly.

\textbf{Optional neutron-density coefficients.} All three segment classes
optionally accept a second, independent set of coefficients for the neutron
density $n_n$ -- same functional form and coordinate as the electron-density
fit above, fitted or supplied separately -- exposing \code{average\_n},
\code{potential\_n}, \pyfunc{residual\_integral\_neutron}, and
\pyfunc{has\_perturbation\_neutron} alongside the electron-only attributes,
required by \pyfunc{evolutor\_perturbative\_segment}'s
\code{include\_matter\_nc=True} (previous section). They default to
\code{None}, reproducing the pre-existing CC-only behaviour exactly; the
layered profile classes that build these segments
(\code{EvenPowerProfileLayered}, \code{PremTabulatedProfile}) load them from
a composition-derived neutron-density table when constructed with
\code{include\_neutron=True} (\cref{ch:earth}), and
\code{AtmospherePolynomialProfile} fits them from a separate neutron-density
sample at the same nodes (\cref{ch:atmosphere}).
\end{implbox}
